% TOP_PROBLEM: {"id": 2512, "title": "Kourovka Notebook Problem 21.3", "category_id": 17, "category": {"id": 17, "name": "group_theory", "display_name": "Group Theory", "description": "Problems about groups, group actions, representations, and related algebraic structures.", "slug": "group-theory", "order_index": 17, "created_at": "2026-07-31T15:26:25.671Z"}, "status": "open", "problem_number": "KOU-21.3", "set_id": 10, "statusQualification": "Part (a) is marked solved in the 30 September2026 primary edition; part (b) remains the retained question."}
% FIRST_POSTED: 2026-10-01
% ENTRY_KIND: statement_only
% UnsolvedMath ID 2512; source catalogue code KOU-21.3
% !TeX program = lualatex
% Definition and sources only; this file is not an LLM solution attempt.
\documentclass[11pt,a4paper]{article}
\usepackage[margin=25mm]{geometry}
\usepackage{fontspec}
\setmainfont{lmroman10-regular.otf}[BoldFont=lmroman10-bold.otf,
 ItalicFont=lmroman10-italic.otf,BoldItalicFont=lmroman10-bolditalic.otf]
\usepackage{amsmath,amssymb,mathtools}
\usepackage{unicode-math}
\setmathfont{latinmodern-math.otf}
\AtBeginDocument{\let\setminus\smallsetminus}
\usepackage{xurl,hyperref}
\hypersetup{colorlinks=true,urlcolor=blue}
\setcounter{secnumdepth}{0}
\setlength{\parindent}{0pt}
\setlength{\parskip}{5pt}
\setlength{\emergencystretch}{3em}
\begin{document}
\section*{Kourovka Notebook Problem 21.3}\label{2512}
{\small UnsolvedMath ID 2512 · KOU-21.3 · Group Theory · Kourovka Notebook - New Problems, Issue 21}
\subsection{Definitions and mathematical statement}
Let $S_n$ be the permutation group of $\{1,\ldots,n\}$ and let $A_n$
be its subgroup of even permutations. For a group $L$, define
$L^{(0)}=L$ and $L^{(j+1)}=[L^{(j)},L^{(j)}]$, where the commutator
subgroup is generated by $[a,b]=a^{-1}b^{-1}ab$. The group is soluble
if $L^{(j)}=1$ for some finite $j$.
For a subgroup $K\le G$ and $x\in G$, put $K^x=x^{-1}Kx$.

The first question asks whether there is an integer $N$ such that for
every $n\ge N$, for either choice $G=A_n$ or $G=S_n$, and every pair
of soluble subgroups $H,K\le G$, there exists $x\in G$ satisfying
\[
 H\cap K^x=\{1\}.
\]
The second asks the sharper assertion with the specific threshold
$N=21$. Thus it includes every $n\ge21$, not only a sufficiently large
subsequence of degrees.

The two soluble subgroups may coincide, and neither is assumed maximal
as a subgroup or maximal among soluble subgroups. The conjugating element
may depend on both subgroups and on the ambient group. For $G=A_n$ it
must be an even permutation; finding a conjugator only in $S_n$ would
not satisfy that instance. The intersection is required to be exactly
the identity, with no nontrivial element common to the two subgroups.
One uniform threshold must handle all allowed pairs simultaneously.

The current primary edition records an affirmative solution of part (a).
Both original parts are defined here; part (b), with the explicit
threshold twenty-one, is the remaining question.
\subsection{Short English statement}
For large symmetric and alternating groups, can any two soluble subgroups be conjugated to intersect trivially, and does degree twenty-one already suffice?
\subsection{Sources}
\begin{itemize}
\item[S1] ULAM.AI and the UnsolvedMath Contributors, supplied problems.json, record 2512 (KOU-21.3), Kourovka Notebook - New Problems, Issue 21. Catalogue identity and source transcription; CC BY 4.0.
\url{https://huggingface.co/datasets/ulamai/UnsolvedMath}.
\item[S2] The Kourovka Notebook, 21st edition, new problems, Problem 21.3. Expanded formulation of the supplied catalogue statement.
\url{https://arxiv.org/pdf/1401.0300v47}.
\end{itemize}
\end{document}
